\subsection{Polynomial cubic alternatives and pullback calculations}
\label{app:application-details}
\label{app:Cartan-source-certificate}

We detail the polynomial potentials of
\cref{sec:cubic-algebraic-existence}, Cartan pullbacks of
\cref{thm:cartan-pullback}, and unequal Lawson pullback of
\cref{ex:lawson-pullback}.

\subsubsection*{The degree-seven defect}
\begin{proof}[Proof of \cref{lem:cubic-SF-defect}]
The product rule gives \eqref{eq:cubic-SF-general-gradient}.
\Cref{lem:cubic-radial-transports} and Euler's identities give
\[
 \begin{aligned}
 \Gamma(G,R)&=14SF,& \Gamma(G,F)&=S^2+4RF^2,\\
 \Gamma(G,S)&=8RSF+12DF,& \Gamma(G,D)&=2S^2F+8RDF.
 \end{aligned}
\]
Differentiating the displayed polynomial \(Q\) in \(R,F,S,D\)
and substituting in \(\cL(G)=Q\Delta G-\Gamma(G,Q)/2\)
proves \eqref{eq:cubic-SF-general-defect}.
\end{proof}

\subsubsection*{The positive decomposition for the polynomial certificate}
Under the hypotheses of \cref{thm:cubic-algebraic-trace-certificate},
the multiplication dictionary and Bochner identity give
\begin{equation}\label{eq:cubic-gradient-identities-only}
 \Delta F=0,\qquad
 \Gamma(F,S)=4RF,\qquad
 \Gamma(S,S)=4RS+12F^2,
\end{equation}
and
\begin{equation}\label{eq:cubic-variable-trace}
 K:=\Delta S=2\tr(L_x^2)\ge16R.
\end{equation}
The trace bound gives $F\not\equiv0$, and
\cref{lem:cubic-unified-bounds} gives
\begin{equation}\label{eq:cubic-gradient-only-S-bound}
 0\le S\le R^2/3.
\end{equation}

Set \(D=F^2\). Then
\(\Gamma(F,D)=2SF\) and \(\Gamma(S,D)=8RD\), so
\cref{lem:cubic-SF-defect} gives
\begin{equation}\label{eq:cubic-variable-trace-gradient}
 Q:=|\nabla G|^2=S^3+12RSF^2+12F^4,
 \qquad \Delta G=(K+8R)F.
\end{equation}
The defect \eqref{eq:cubic-SF-general-defect} specializes to
\begin{align}
 \frac{\cL(G)}F
 &=(K-16R)S^3+F^2S(12RK-126S)
                       +12F^4(K-6R)\nonumber\\
 &=(K-16R)Q+150R^2SF^2
       +126F^2S\left(\frac{R^2}3-S\right)+120RF^4.
 \label{eq:cubic-variable-trace-positive-decomposition}
\end{align}
The quotient is understood by polynomial extension across $F=0$.
Every term on the last line is nonnegative by
\eqref{eq:cubic-variable-trace} and
\eqref{eq:cubic-gradient-only-S-bound}. If $F\ne0$, Euler's
identity implies $S>0$ and $R>0$, so the term $150R^2SF^2$
is strictly positive. Thus $F\cL(G)>0$ away from $F=0$.

The geometric comparison is proved in
\cref{thm:cubic-algebraic-trace-certificate}.


\paragraph{The Cartan source certificate.}

For \(F,Q,G\) in \cref{thm:cartan-pullback}, the pullback identities give
\begin{align}
 |\nabla G|^2&=4RQ(9Q^3+16F^2),
 \label{eq:Cartan-pullback-gradient}\\
 \cL(G)&=8QF\Bigl[
 9\bigl((4m-48)R^2-5Q\bigr)Q^3\notag\\
 &\hspace{31mm}
 +16\bigl((4m+2)R^2-5Q\bigr)F^2\Bigr].
 \label{eq:Cartan-pullback-defect}
\end{align}
Since \(Q\leq R^2\), the bracket is positive when \(Q>0\).
Thus \(G\cL(G)\geq0\) throughout the regular set.

\paragraph{The unequal Lawson pullback.}
For the target data $T,f,H$ in \cref{ex:lawson-pullback},
\[
 \Delta f=2,\qquad |\nabla f|^2=80T+4f,\qquad
 \Gamma(T,f)=4f.
\]
In particular $\cL(f)=-160f$, so its regular zero set is minimal.
The product rule gives
\begin{align*}
 |\nabla H|^2&=4T(20T^2+Tf+3f^2),\\
 \cL(H)&=8Tf(60T^2+7Tf+33f^2),\\
 2T\cL(H)-4H|\nabla H|^2
 &=32T^2f(20T^2+3Tf+15f^2).
\end{align*}
The last quadratic factor is positive for $T>0$: its discriminant
as a polynomial in $f$ is $-1191T^2$ and its leading coefficient
is positive. Hence \eqref{eq:pullback-target-sign} holds with
$\ell=4$ and $\kappa=1/2$, strictly off $f=0$.
For the source weights $(22,5,6)$, one has $N=64$ and $h=11$.
Substitution in \eqref{eq:bihomogeneous-Hardy} gives
\[
 Q_C(u)\geq\frac{2273}{4}\int_Cr^{-2}u^2,
 \qquad u\in C_c^\infty(C_{\reg}\setminus\{0\}),
 \quad C=\{F=0\},
\]
since $(64-17)^2/4+8\cdot11-72=2273/4$.
